\section{Implementación numérica}

\subsection{Introducción}

\hspace{1em}El código empleado en este trabajo es MSTM (\textit{Multiple 
Sphere T-Matrix}), una implementación en Fortran desarrollada por Mackowski 
\cite{Mackowski_code} del formalismo de la T-matrix para el scattering 
electromagnético por agregados de esferas. El código resuelve de forma 
autoconsistente el sistema de Foldy-Lax (\ref{eq:foldy-lax}), calcula la 
T-matrix del agregado y proporciona las magnitudes observables: secciones 
eficaces, matriz de Mueller y campo cercano.

Una característica central del código es que trabaja completamente en 
unidades adimensionales. La longitud de onda de la radiación incidente no 
aparece de forma explícita: todas las longitudes se expresan en unidades del 
número de onda $k = 2\pi/\lambda$ del medio circundante. En consecuencia, 
los radios de las esferas se introducen directamente como parámetros de 
tamaño $x_j = k a_j$, y las posiciones de los centros se expresan en las 
mismas unidades. Esta convención implica que un cambio en la longitud de onda 
es equivalente a un reescalado simultáneo de todas las dimensiones geométricas 
del sistema. Cuando se desea trabajar con unidades físicas, el código 
proporciona el parámetro \texttt{length\_scale\_factor}, que reescala 
internamente todas las posiciones y radios antes de comenzar el cálculo.


\subsection{Estructura general del programa}

\subsubsection{Flujo de ejecución}

\hspace{1em}El funcionamiento del programa sigue directamente la cadena 
lógica del formalismo de la T-matrix descrito en la sección anterior:

\begin{enumerate}
  \item Lectura de los parámetros de entrada: parámetros de tamaño 
        $x_j = ka_j$, posiciones $\{\mathbf{r}_j\}$ e índices de refracción 
        complejos $\{m_j\}$ de las esferas.
  \item Inicialización de la geometría del sistema: cálculo de los vectores 
        de separación $\mathbf{d}_{jl} = \mathbf{r}_j - \mathbf{r}_l$ y 
        determinación del radio de circunscripción del agregado.
  \item Cálculo de los coeficientes de Mie para cada esfera, que determinan 
        su T-matrix individual $\mathbb{T}_j$ (\ref{eq:miecoeffs}). El orden 
        multipolar $n_{max}$ se determina automáticamente mediante el criterio 
        de Wiscombe $n_{max} \approx x + 4x^{1/3} + 5$. Aunque en este 
        trabajo todas las esferas son idénticas ($\mathbb{T}_j = \mathbb{T}$), 
        el código calcula la T-matrix de forma independiente para cada partícula.
  \item Cálculo de los coeficientes de traslación $\mathbb{H}_{jl}$, que 
        permiten expresar el campo dispersado por la esfera $l$ en el sistema 
        de referencia de la esfera $j$.
  \item Construcción y resolución del sistema lineal de Foldy-Lax 
        (\ref{eq:foldy-lax}), que proporciona los coeficientes multipolares 
        $\{\mathbf{p}_j\}$ de forma autoconsistente para dos polarizaciones 
        ortogonales de la onda incidente de forma simultánea.
  \item Construcción de la T-matrix del agregado $\mathbb{T}^{cluster}$ y 
        cálculo de los observables: secciones eficaces, elementos de la 
        matriz de Mueller y parámetro de asimetría.
  \item Evaluación opcional del campo electromagnético en la región cercana 
        al sistema de esferas.
\end{enumerate}

\subsubsection{Organización en módulos}

\hspace{1em}El código está organizado en módulos, cada uno responsable de 
una etapa del cálculo:

\begin{itemize}
  \item \texttt{numconstants}: constantes matemáticas y coeficientes 
        precalculados que son reutilizados a lo largo del programa, como 
        los coeficientes binomiales $\binom{n+l}{n}^{1/2}$ y los coeficientes 
        de las funciones vectoriales de onda esférica.
  \item \texttt{specialfuncs}: implementación de las funciones especiales 
        necesarias: funciones de Riccati-Bessel $\psi_n$ y Riccati-Hankel 
        $\xi_n$, polinomios de Legendre asociados, funciones angulares 
        $\pi_n$ y $\tau_n$, y coeficientes de traslación $\mathbb{H}_{jl}$.
  \item \texttt{spheredata}: almacenamiento y gestión de la geometría del 
        sistema: posiciones, radios, índices de refracción y jerarquía de 
        esferas anidadas.
  \item \texttt{mie}: cálculo de los coeficientes de Mie $a_n$, $b_n$ y 
        construcción de la T-matrix individual de cada esfera, con soporte 
        para medios ópticamente activos.
  \item \texttt{translation}: implementación del teorema de adición de las 
        VSWF y cálculo de los operadores de traslación $\mathbb{H}_{jl}$ 
        entre pares de esferas.
  \item \texttt{solver}: construcción y resolución numérica del sistema de 
        Foldy-Lax, mediante factorización LU o método iterativo de gradiente 
        biconjugado.
  \item \texttt{scatprops}: cálculo de los observables a partir de la 
        T-matrix del agregado: secciones eficaces, elementos de la matriz 
        de Mueller, parámetro de asimetría $g$ y albedo $\omega_0$.
  \item \texttt{nearfield}: evaluación del campo electromagnético total en 
        una malla espacial en la región cercana al sistema de esferas.
  \item \texttt{inputinterface}: lectura de parámetros de entrada, control 
        del flujo de ejecución y escritura de resultados.
\end{itemize}


\subsection{Métodos numéricos empleados}

\subsubsection{Evaluación de funciones especiales y orden multipolar}

\hspace{1em}La construcción de los coeficientes de Mie (\ref{eq:miecoeffs}), 
de las funciones angulares $\pi_n$ y $\tau_n$ (\ref{eq:S1S2}) y de los 
coeficientes de traslación $\mathbb{H}_{jl}$ requiere evaluar las funciones 
de Riccati-Bessel $\psi_n$ y Riccati-Hankel $\xi_n$ y sus derivadas. Estas 
derivadas se calculan a partir de las relaciones de recurrencia analíticas 
que satisfacen, como

\begin{equation}
  \frac{d}{d\rho}\left[\rho\,j_n(\rho)\right] 
  = \rho\,j_{n-1}(\rho) - n\,j_n(\rho),
\end{equation}

y relaciones análogas para las funciones de Hankel. El uso de recurrencias 
en lugar de aproximaciones por diferencias finitas garantiza la estabilidad 
numérica del cálculo incluso para órdenes elevados. De forma análoga, las 
funciones angulares $\pi_n$ y $\tau_n$ se evalúan mediante sus relaciones 
de recurrencia a partir de los valores iniciales $\pi_0 = 0$, $\pi_1 = 1$.

El orden multipolar máximo $n_{max}$ retenido en la expansión se determina 
automáticamente para cada esfera mediante el criterio de Wiscombe:

\begin{equation}
  n_{max} \approx x + 4\,x^{1/3} + 5,
\end{equation}

que garantiza la convergencia de la serie para cualquier valor del parámetro 
de tamaño $x$. La dimensión del sistema de Foldy-Lax es entonces 
$N \times 2n_{max}(n_{max}+2)$ por esfera.


\subsubsection{Resolución del sistema de Foldy-Lax}

\hspace{1em}El sistema de Foldy-Lax (\ref{eq:foldy-lax}) adopta la forma 
de un sistema lineal

\begin{equation}
  \mathbb{M}_{sys}\,\mathbf{p} = \mathbf{b},
\end{equation}

donde $\mathbf{p}$ contiene los coeficientes multipolares de todos los 
campos dispersados, $\mathbf{b}$ recoge las contribuciones del campo 
incidente sobre cada esfera, y $\mathbb{M}_{sys}$ incorpora la respuesta 
individual de cada esfera a través de su T-matrix y el acoplamiento entre 
partículas a través de los operadores de traslación $\mathbb{H}_{jl}$. La 
matriz es compleja, densa y de dimensión $N \times 2n_{max}(n_{max}+2)$.

El código implementa dos métodos de resolución, seleccionables por el 
usuario. Para sistemas de tamaño moderado se emplea la \textit{factorización 
LU}, que descompone la matriz como

\begin{equation}
  P\,\mathbb{M}_{sys} = L\,U,
\end{equation}

donde $P$ es una matriz de permutación asociada al pivotado parcial, $L$ es 
triangular inferior con unos en la diagonal y $U$ es triangular superior. 
La descomposición se obtiene mediante eliminación gaussiana,

\begin{equation}
  a_{ij}^{(k+1)} = a_{ij}^{(k)} - 
  \frac{a_{ik}^{(k)}}{a_{kk}^{(k)}}\,a_{kj}^{(k)}, 
  \quad i > k,\; j \geq k,
\end{equation}

y el sistema se resuelve en dos etapas de sustitución hacia adelante y hacia 
atrás. El coste computacional es $\mathcal{O}(N^3)$.

Para sistemas de mayor tamaño se emplea el método iterativo del 
\textit{gradiente biconjugado} (BiCG), que no requiere construir 
explícitamente $\mathbb{M}_{sys}$, sino únicamente evaluar productos 
matriz-vector $\mathbb{M}_{sys}\,\mathbf{v}$, calculados directamente 
a partir de las interacciones entre esferas mediante los operadores de 
traslación. La solución se actualiza iterativamente como

\begin{equation}
  \mathbf{p}_{k+1} = \mathbf{p}_k + \alpha_k\,\mathbf{d}_k,
\end{equation}

donde $\mathbf{d}_k$ es una dirección de búsqueda y $\alpha_k$ minimiza 
el residuo en cada paso, hasta satisfacer el criterio de convergencia

\begin{equation}
  \frac{\|\mathbf{r}_k\|}{\|\mathbf{b}\|} < \varepsilon,
\end{equation}

con un coste por iteración de $\mathcal{O}(N)$.

El sistema se resuelve simultáneamente para dos polarizaciones ortogonales 
de la onda incidente, de modo que en una sola ejecución se obtienen los 
coeficientes $\{\mathbf{p}_j\}$ correspondientes a ambas polarizaciones, 
lo que permite construir directamente todos los elementos de la matriz de 
Mueller.


\subsection{Cálculo del campo cercano}

\subsubsection{Definición}

\hspace{1em}Una vez resuelto el sistema de Foldy-Lax, los coeficientes 
$\{\mathbf{p}_j\}$ permiten evaluar el campo electromagnético en cualquier 
punto del espacio. El campo total se descompone como

\begin{equation}
  \mathbf{E}(\mathbf{r}) = \mathbf{E}^{inc}(\mathbf{r}) 
  + \sum_{j=1}^{N} \mathbf{E}^{sca}_j(\mathbf{r}),
\end{equation}

donde cada contribución dispersada se evalúa mediante la expansión en VSWF 
centrada en $\mathbf{r}_j$. A diferencia del campo lejano, donde se toma 
el límite $kr \to \infty$ y las funciones de Hankel se aproximan por ondas 
esféricas salientes, en el campo cercano las funciones de Hankel esféricas 
deben evaluarse en su forma completa, sin ninguna aproximación asintótica. 
Esto permite analizar la distribución espacial local de la energía 
electromagnética en el entorno inmediato de las esferas, incluyendo 
fenómenos como la concentración de campo en las zonas de contacto entre 
partículas adyacentes.

\subsubsection{Implementación en el código}

\hspace{1em}El cálculo del campo cercano se lleva a cabo mediante el módulo 
\texttt{nearfield}, que evalúa los campos eléctrico y magnético sobre una 
malla espacial cartesiana definida por el usuario. Para cada punto 
$\mathbf{r} = (x,y,z)$ de la malla, el programa calcula la posición 
relativa respecto al centro de cada esfera, $\mathbf{r} - \mathbf{r}_j$, 
evalúa las VSWF en esa posición y suma las contribuciones de todos los 
modos multipolares y todas las esferas, añadiendo la contribución del campo 
incidente para obtener el campo total. El proceso se realiza de forma 
independiente para las dos polarizaciones ortogonales, de modo que la salida 
contiene, para cada punto de la malla, las tres componentes complejas de los 
campos eléctrico y magnético para cada polarización.


\subsection{Interpretación de resultados}

\hspace{1em}Para luz no polarizada, los observables se obtienen promediando 
las contribuciones de las dos polarizaciones ortogonales calculadas. Dado 
que estas polarizaciones son incoherentes entre sí, la combinación se 
realiza sobre magnitudes cuadráticas. Para el campo cercano, la intensidad 
media se expresa como

\begin{equation}
  \langle |\mathbf{E}|^2 \rangle = 
  \frac{|\mathbf{E}_\parallel|^2 + |\mathbf{E}_\perp|^2}{2}.
\end{equation}

Para las magnitudes de campo lejano, el elemento $S_{11}(\theta)$ de la 
matriz de Mueller obtenido directamente del código corresponde ya al 
promedio sobre ambas polarizaciones, mientras que los demás elementos 
$S_{12}$, $S_{33}$ y $S_{34}$ permiten caracterizar el estado de 
polarización del campo dispersado a través del grado de polarización lineal 
$P_L(\theta) = -S_{12}(\theta)/S_{11}(\theta)$.